% Options for packages loaded elsewhere
\PassOptionsToPackage{unicode}{hyperref}
\PassOptionsToPackage{hyphens}{url}
\PassOptionsToPackage{dvipsnames,svgnames,x11names}{xcolor}
\documentclass[
  10pt,
  fontset=fandol]{ctexart}
\usepackage{xcolor}
\usepackage[margin=16mm]{geometry}
\usepackage{amsmath,amssymb}
\setcounter{secnumdepth}{-\maxdimen} % remove section numbering
\usepackage{iftex}
\ifPDFTeX
  \usepackage[T1]{fontenc}
  \usepackage[utf8]{inputenc}
  \usepackage{textcomp} % provide euro and other symbols
\else % if luatex or xetex
  \usepackage{unicode-math} % this also loads fontspec
  \defaultfontfeatures{Scale=MatchLowercase}
  \defaultfontfeatures[\rmfamily]{Ligatures=TeX,Scale=1}
\fi
\usepackage{lmodern}
\ifPDFTeX\else
  % xetex/luatex font selection
\fi
% Use upquote if available, for straight quotes in verbatim environments
\IfFileExists{upquote.sty}{\usepackage{upquote}}{}
\IfFileExists{microtype.sty}{% use microtype if available
  \usepackage[]{microtype}
  \UseMicrotypeSet[protrusion]{basicmath} % disable protrusion for tt fonts
}{}
\makeatletter
\@ifundefined{KOMAClassName}{% if non-KOMA class
  \IfFileExists{parskip.sty}{%
    \usepackage{parskip}
  }{% else
    \setlength{\parindent}{0pt}
    \setlength{\parskip}{6pt plus 2pt minus 1pt}}
}{% if KOMA class
  \KOMAoptions{parskip=half}}
\makeatother
\setlength{\emergencystretch}{3em} % prevent overfull lines
\providecommand{\tightlist}{%
  \setlength{\itemsep}{0pt}\setlength{\parskip}{0pt}}
\usepackage{amsmath,amssymb,mathtools,needspace}
\xeCJKDeclareCharClass{CJK}{"2032,"0400->"04FF,"2160->"216B,"2460->"2473,"25A0,"25C7,"25CB,"25CF}
\IfFontExistsTF{Arial Unicode MS}{\xeCJKsetup{AutoFallBack=true}\setCJKfallbackfamilyfont{\CJKrmdefault}{Arial Unicode MS}}{}
\setcounter{secnumdepth}{0}
\setlength{\emergencystretch}{2em}
\allowdisplaybreaks
\usepackage{bookmark}
\IfFileExists{xurl.sty}{\usepackage{xurl}}{} % add URL line breaks if available
\urlstyle{same}
\hypersetup{
  pdftitle={二分手续},
  colorlinks=true,
  linkcolor={Maroon},
  filecolor={Maroon},
  citecolor={Blue},
  urlcolor={Blue},
  pdfcreator={LaTeX via pandoc}}

\title{二分手续}
\author{}
\date{}

\begin{document}
\maketitle

\subsubsection{11.2.2　二分手续}\label{ux4e8cux5206ux624bux7eed}

为使并行算法设计的原理与方法简单而和谐，这里推荐一种设计技术------二分技术。

\textbf{二分技术的设计原理是，反复地将所给计算问题加工成规模减半的同类问题，直到规模足够小（通常当规模为 1）时直接得出问题的解。}

\textbf{需要强调的是，与倍增技术不同，二分技术不是着眼于问题的分裂，而是立足于问题的加工。今后所说的二分手续，是指将问题规模减半的加工手续。}

譬如，对于 \(N\) 项和式

\[
S=\sum_{i=0}^{N-1}a_i,
\]

\href{../../source-images/pdf-281.jpeg}{原页图像}

若将其前后对应项两两合并，即可加工成一个规模减半的 \(N_1=N/2\) 项和式

\[
S=(a_0+a_{N-1})+(a_1+a_{N-2})+\cdots+(a_{N/2-1}+a_{N/2}).
\]

这种二分手续联系着大数学家 Gauss 幼年时代的一个小故事。

有一天，算术课老师要小学生们计算前 100 个自然数的和 \(S=1+2+\cdots+99+100\)。当班上其他同学忙于逐项累加而头昏脑胀的时候，小 Gauss 却机智地发现，所给和式前后对应项的和均等于 101，因而所求和值为 \(S=101\times50=5\,050\)。这种简捷的快速算法可以视作二分手续的巧妙应用。

\begin{center}\rule{0.5\linewidth}{0.5pt}\end{center}

\subsection{相关原页校注（agent 补充）}\label{ux76f8ux5173ux539fux9875ux6821ux6ce8agent-ux8865ux5145}

以下校注来自本节覆盖的原页，可能也涉及同页相邻小节；原文照录与校正说明分开保留。

\subsubsection{原书 PDF 页 281 的校注}\label{ux539fux4e66-pdf-ux9875-281-ux7684ux6821ux6ce8}

\href{../pages/pdf-281.md}{完整原页转录} · \href{../../source-images/pdf-281.jpeg}{原页图像}

校注记录：CH11-E001。

\begin{quote}
校注（agent 补充，CH11-E001）：本页两处回引均印作``（10.2.1）''，已照录；从本节的数列求和定义及上下文判断，应回引本章的式（11.2.1）。这是原书交叉引用疑误，不是 OCR 改写。
\end{quote}

\subsection{本节覆盖原页的校注记录}\label{ux672cux8282ux8986ux76d6ux539fux9875ux7684ux6821ux6ce8ux8bb0ux5f55}

下列记录按原页关联，含同页相邻小节的校注；计算时同时读取上文校注和完整原页。

\begin{itemize}
\tightlist
\item
  \texttt{CH11-E001}：原书 PDF 页 281
\end{itemize}

\end{document}
